\begin{ZhChapter}

\chapter{CBR-APIKGAgent 系統架構與方法}

本章說明本研究所提出之 CBR-APIKGAgent 的系統架構與核心方法。承接前一章對多步驟 API 任務、經驗重用與動態計畫修復相關研究之整理，本章將重點放在兩項改良機制：其一為由訓練資料成功任務軌跡建立的原子化經驗重用機制，其二為針對執行失敗進行局部計畫調整的動態局部修復機制。以下首先介紹整體系統架構與任務流程，接著分別說明原子化經驗的表示、檢索與注入方式，以及局部修復模組如何根據失敗資訊進行插入、替換或刪除等修復操作。

\section{CBR-APIKGAgent 系統總覽}

CBR-APIKGAgent 的整體架構可視為在 APIKGAgent 原有流程中加入兩個輔助層：其一為離線原子化經驗庫與語意檢索流程，其二為執行失敗後的局部修復流程。以下先以整體架構圖呈現任務從初始規劃、步驟執行、局部修復到重新規劃的資料流與控制流，再整理主要元件之功能與相互關係。

\subsection{整體架構與主要元件}

CBR-APIKGAgent 以 APIKGAgent 為基礎，延續其 API 知識圖譜、API 搜尋、任務規劃、步驟式執行與重新規劃流程，並進一步加入原子化經驗重用與動態局部修復機制。整體而言，本研究並未改變 APIKGAgent 以「先搜尋 API、再產生計畫、最後逐步執行」為核心的任務執行模式，而是在規劃與錯誤處理階段加入案例式經驗檢索，使代理人在面對相似任務結構或局部執行失敗時，能取得來自訓練資料成功任務軌跡的策略參考。

CBR-APIKGAgent 的整體架構如圖 \ref{fig:cbr_apikgagent_architecture} 所示。系統以使用者任務為輸入，先透過 API 知識圖譜搜尋候選 API，並由任務規劃器產生初始任務計畫。任務執行期間，任務執行器會逐步執行計畫並保存中間結果；當步驟失敗時，系統會先嘗試由局部修復模組修正局部計畫，若局部修復無法處理，則回退至重新規劃器產生後續計畫。原子化經驗檢索器則同時支援初始規劃、重新規劃與局部修復三個決策階段，使代理人能取得與當前任務或錯誤情境相關的策略參考。

\begin{figure}[htbp]
    \centering
    \IfFileExists{figures/chapter3/cbr-apikgagent-architecture.png}{%
        \includegraphics[
            width=\textwidth,
            height=0.78\textheight,
            keepaspectratio
        ]{figures/chapter3/cbr-apikgagent-architecture.png}%
    }{%
        \fbox{\begin{minipage}{0.92\textwidth}
            \centering
            \textbf{CBR-APIKGAgent 整體架構示意圖}\\[0.8em]
            \begin{tabular}{p{0.24\textwidth} p{0.62\textwidth}}
                \hline
                \textbf{離線經驗層} &
                訓練資料成功任務軌跡
                $\rightarrow$
                原子化經驗庫
                $\rightarrow$
                規劃導向檢索／修復導向檢索 \\
                \hline
                \textbf{任務主流程} &
                使用者任務
                $\rightarrow$
                API 知識圖譜搜尋
                $\rightarrow$
                任務規劃器產生初始計畫
                $\rightarrow$
                任務執行器逐步執行 \\
                \hline
                \textbf{成功路徑} &
                步驟成功
                $\rightarrow$
                保存中間結果
                $\rightarrow$
                執行下一步
                $\rightarrow$
                完成任務 \\
                \hline
                \textbf{失敗處理} &
                步驟失敗
                $\rightarrow$
                形成錯誤證據
                $\rightarrow$
                局部修復模組（插入／替換／刪除）
                $\rightarrow$
                修復後回到任務執行器 \\
                \hline
                \textbf{回退機制} &
                局部修復失敗或超出範圍
                $\rightarrow$
                重新規劃器產生後續計畫
                $\rightarrow$
                回到任務執行器 \\
                \hline
            \end{tabular}\\[0.8em]
            \hrule
            \vspace{0.6em}
            圖中重點：原子化經驗檢索支援初始規劃、局部修復與重新規劃；局部修復優先處理可局部調整之錯誤，必要時再回退至重新規劃。
        \end{minipage}}%
    }
    \caption{CBR-APIKGAgent 整體架構}
    \label{fig:cbr_apikgagent_architecture}
\end{figure}

圖 \ref{fig:cbr_apikgagent_architecture} 中，案例式經驗檢索與動態局部修復雖然在任務失敗時會互相配合，但兩者在系統中扮演不同角色。案例式經驗檢索負責提供與當前任務或錯誤情境相關的策略參考，屬於決策輔助機制；動態局部修復則負責根據失敗資訊修改局部計畫，屬於錯誤處理機制。

CBR-APIKGAgent 的主要元件可分為六個部分：API 知識圖譜與 API 搜尋模組、任務規劃器、任務執行器、原子化經驗檢索器、動態局部修復模組，以及重新規劃器。為補充圖中各模組之功能，表 \ref{tab:cbr_apikgagent_components} 整理各元件在系統中的責任。

\begin{table}[htbp]
    \centering
    \caption{CBR-APIKGAgent 主要元件與功能}
    \label{tab:cbr_apikgagent_components}
    \small
    \renewcommand{\arraystretch}{1.15}
    \begin{tabularx}{0.95\textwidth}{
        >{\raggedright\arraybackslash}p{0.23\textwidth}
        >{\raggedright\arraybackslash}p{0.16\textwidth}
        >{\raggedright\arraybackslash}X
    }
        \hline
        \textbf{元件} & \textbf{角色} & \textbf{功能說明} \\
        \hline
        API 知識圖譜／搜尋模組 &
        API 資訊來源 &
        依據任務目標或修復需求搜尋相關 API，提供規劃與修復所需資訊。 \\[0.35em]
        任務規劃器 &
        初始規劃 &
        根據使用者任務、候選 API 與規劃經驗產生初始任務計畫。 \\[0.35em]
        任務執行器 &
        步驟執行 &
        依照任務計畫逐步執行 API 呼叫與資料處理，並保存中間結果與失敗資訊。 \\[0.35em]
        原子化經驗檢索器 &
        策略參考 &
        從離線經驗庫檢索相關經驗，提供初始規劃、重新規劃與局部修復作為策略參考。 \\[0.35em]
        局部修復模組 &
        錯誤修復 &
        在步驟失敗時，根據錯誤證據與修復經驗，以插入、替換或刪除步驟調整局部計畫。 \\[0.35em]
        重新規劃器 &
        回退處理 &
        當局部修復無法處理錯誤時，根據任務目標、既有計畫與失敗資訊重新產生後續計畫。 \\
        \hline
    \end{tabularx}
\end{table}

在資料流設計上，CBR-APIKGAgent 以任務狀態作為各階段之間共享上下文的載體。任務狀態保存任務目標、候選 API、目前計畫、已完成步驟結果與失敗資訊，使規劃、執行、局部修復與重新規劃能在同一脈絡下進行。較細部的狀態欄位、資料保存方式與程式傳遞流程，則於第四章系統實作中說明。

\subsection{架構擴充重點}

本研究以 APIKGAgent 作為基礎代理人架構，延續其以 API 知識圖譜輔助 API 搜尋、任務規劃、步驟式執行與重新規劃的基本流程。在此基礎上，CBR-APIKGAgent 並未改變原有的任務執行主軸，而是針對多步驟 API 任務中經驗重用不足與局部錯誤不易修復兩項問題進行擴充。

第一項擴充為原子化經驗重用機制。系統將訓練資料中的成功任務軌跡轉換為可檢索的經驗單元，並於初始規劃、重新規劃與局部修復階段提供與當前任務或錯誤情境相關的策略參考。此設計的目的不是直接複製過去任務軌跡，而是使代理人在面對相似資料依賴、規劃條件或常見陷阱時，能取得可供決策參考的經驗知識。

第二項擴充為動態局部修復機制。當任務執行過程中發生步驟失敗時，系統不立即進入整體重新規劃，而是先根據錯誤證據、相依資料、候選修復 API 與修復經驗評估是否存在可行的局部修復方式。若系統能產生局部調整方案，則嘗試透過插入、替換或刪除步驟修正局部計畫；若無法形成可行修復方案、修復次數達到限制，或錯誤訊號不足以支持局部修復，則回退至重新規劃。

此外，CBR-APIKGAgent 正式使用的經驗庫由訓練資料中的成功任務軌跡離線建立，測試期間不將新產生的修復紀錄寫回正式檢索索引。因此，本研究方法屬於離線經驗庫輔助的案例式修復架構，而非測試期間更新正式經驗庫的方法。

\subsection{任務執行與錯誤修復流程}

CBR-APIKGAgent 的任務執行流程可分為初始規劃、逐步執行、局部修復與重新規劃四個階段。各階段之功能如下：

\begin{enumerate}
    \item \textbf{初始規劃：}系統接收使用者任務後，透過 API 知識圖譜搜尋與任務相關的候選 API，並從原子化經驗庫中檢索規劃相關經驗。任務規劃器根據任務目標、候選 API 與檢索經驗產生初始任務計畫。

    \item \textbf{逐步執行：}產生初始計畫後，任務執行器依序執行計畫中的 API 呼叫與資料處理步驟。成功步驟的輸出會保存於任務狀態中，作為後續步驟的相依資料；若所有步驟皆成功完成，系統即依任務需求提交最終答案或完成環境狀態變更。

    \item \textbf{局部修復：}若某一步驟執行失敗，系統會收集失敗步驟、錯誤訊息、相關中間結果與 API 執行回饋，形成局部修復所需的錯誤證據。局部修復模組會結合錯誤證據、候選修復 API 與修復相關經驗，嘗試產生可行的局部修復方案。若修復成功，系統更新局部計畫並回到任務執行器繼續執行。

    \item \textbf{重新規劃：}若局部修復模組無法選出適合的修復 API、無法產生可執行的局部修復方案，或局部修復次數達到限制，系統則回退至重新規劃器。重新規劃器會參考原始任務、既有計畫、失敗資訊、已完成步驟結果與檢索經驗，重新產生後續計畫，並將修正後的計畫交回任務執行器繼續執行。
\end{enumerate}

整體而言，CBR-APIKGAgent 採取「先局部修復、必要時重新規劃」的錯誤處理流程。此流程的設計目的並非取代重新規劃，而是讓系統在錯誤僅侷限於局部步驟時，能先嘗試保留原本仍有效的計畫與資料流程；當局部修復不足以處理錯誤時，再回退至重新規劃。透過原子化經驗檢索與動態局部修復的結合，本研究進一步探討案例式策略參考是否有助於改善多步驟 API 任務中的規劃品質與錯誤修復表現。

\section{原子化經驗重用機制}

原子化經驗重用機制的目的，是將訓練資料中成功任務軌跡所隱含的規劃與資料流策略，轉換為可於不同任務間重用的經驗知識。相較於直接保存完整任務軌跡，本研究將經驗拆解為較小的策略單元，使其能作為初始規劃、重新規劃與局部修復階段的參考，而非要求代理人複製過去任務的完整步驟。

\subsection{經驗來源與使用範圍}

本研究使用的正式經驗庫由訓練資料中的成功任務軌跡離線建立，測試期間不新增經驗至正式經驗庫，也不更新正式檢索索引。此設計使經驗來源與實驗評估資料保持切割，避免測試期間產生的執行紀錄影響後續任務評估。

雖然系統在局部修復過程中可能產生修復紀錄，但本研究並未將此類紀錄納入正式經驗庫。主要原因在於，單一次局部修復紀錄不一定代表整體任務成功，且該修復是否正確需依後續資料流與最終任務結果判斷。若未經驗證即將修復紀錄轉換為可重用經驗，可能使錯誤或不完整的修復策略影響後續規劃。因此，本研究將測試期間的修復紀錄保留為分析資料，而非正式檢索經驗的一部分。

在經驗萃取階段，本研究以訓練資料中已知成功的任務為來源，參考任務規格、公開資料、標準解答、評估條件與實際 API 呼叫軌跡等資訊，萃取其中可跨任務重用的規劃原則與資料依賴模式。由於本研究關注的是多步驟 API 任務中的規劃與修復能力，因此經驗萃取時不以保留完整 API 呼叫序列為目標，而是著重於任務成功過程中可泛化的策略知識。

\subsection{原子化經驗定義}

本研究將原子化經驗定義為一個可跨任務重用的最小策略單元，用以描述特定任務情境下應保留的規劃原則、資料流模式、必要檢查與常見陷阱。換言之，案例式學習提供本研究的經驗重用觀點，而原子化經驗則是本研究用來表示與檢索案例知識的具體設計。

為避免與完整案例、API 模板或可執行技能混淆，原子化經驗在本研究中的定位如下：

\begin{itemize}
    \item \textbf{不同於完整案例}：原子化經驗不重播完整任務軌跡，而是抽取其中可跨任務重用的策略片段。
    \item \textbf{不同於 API 模板}：原子化經驗不保存固定 API 呼叫序列或參數值，而是描述資料流、檢查條件與決策原則。
    \item \textbf{不同於技能模組}：原子化經驗不是可直接執行的技能（skill），而是提供給規劃與修復流程參考的策略知識。
    \item \textbf{不同於線上長期記憶}：原子化經驗庫由訓練資料中的成功任務軌跡離線建立，測試期間不將新紀錄寫回正式經驗庫。
\end{itemize}

因此，每筆原子化經驗主要描述下列內容：

\begin{itemize}
    \item \textbf{適用情境}：何種任務或錯誤情境下可考慮使用此經驗。
    \item \textbf{核心原則}：當前任務應保留的主要規劃原則。
    \item \textbf{資料依賴}：需要維持的資料流、候選集合或目標集合關係。
    \item \textbf{常見陷阱}：應避免的規劃或修復方式。
\end{itemize}

具體而言，一筆原子化經驗通常對應到一類常見規劃需求或錯誤風險，例如完整取得多頁資料、從原始資料集合中建立正確目標集合、驗證候選結果是否符合任務來源、在寫入前比較目前狀態與目標狀態，或維持輸出格式契約等。這些經驗並不假設當前任務與來源任務完全相同，而是在任務結構、資料依賴或錯誤情境相似時，為代理人提供可參考的規劃與修復原則。

例如，若訓練任務要求代理人對某一條件下的多筆資料進行更新，原子化經驗不會保存「呼叫某個特定 API 並傳入某個固定參數」這類具體操作，而是萃取出「在執行寫入動作前，必須先由來源集合建立符合條件的目標集合，並確認每個目標項目的目前狀態」這類可跨任務重用的規劃原則。

綜合上述定義，原子化經驗在本研究中具有三項特性：

\begin{itemize}
    \item \textbf{粒度較小}：經驗粒度應小於完整任務軌跡，避免將來源任務的特定步驟或 API 名稱直接搬移到新任務。
    \item \textbf{內容可泛化}：經驗內容應描述可泛化的資料流與決策原則，例如目標集合建構、候選驗證、批次完整性與寫入保護。
    % TODO-FINAL-VERSION: 下列「適用性評估、規劃約束」偏向 v2／修正版經驗選擇設計。若最後採 v4，可改為較中性的「經檢索後作為策略參考或修復引導」；若採 v2，則保留並於第四章補上評分與約束注入細節。
    \item \textbf{作為策略約束與修復引導}：原子化經驗經適用性評估後，可轉換為規劃約束或修復引導，用以輔助計畫內容的檢查與調整。實際計畫與修復方案仍需依據當前任務、候選 API、錯誤資訊與執行狀態產生。
\end{itemize}

\subsection{原子化經驗表示方式}

為了使經驗能被檢索、排序並注入不同決策階段，本研究將每筆原子化經驗表示為結構化知識。表 \ref{tab:atomic_experience_schema} 中的欄位用以說明本研究如何將成功任務軌跡抽象化為可檢索經驗；實際資料格式、儲存方式與載入流程則於第四章系統實作中說明。

\begin{table}[htbp]
    \centering
    \caption{原子化經驗主要欄位與用途}
    \label{tab:atomic_experience_schema}
    \small
    \renewcommand{\arraystretch}{1.15}
    \begin{tabularx}{0.95\textwidth}{
        >{\raggedright\arraybackslash}p{0.25\textwidth}
        >{\raggedright\arraybackslash}X
    }
        \hline
        \textbf{欄位} & \textbf{用途說明} \\
        \hline
        \texttt{experience\_type} &
        經驗類型標籤，用於描述此經驗所屬的粗略策略類別，例如目標集合建構、候選驗證、批次完整性或輸出格式契約。 \\[0.35em]
        \texttt{source\_task\_id} &
        來源訓練任務識別碼，用於追溯經驗來源與後續分析。 \\[0.35em]
        \texttt{source\_task\_goal} &
        來源訓練任務目標，用於語意檢索與理解經驗適用背景。 \\[0.35em]
        \texttt{when\_to\_use} &
        描述此經驗適用的任務情境或錯誤情境。 \\[0.35em]
        \texttt{core\_principle} &
        此經驗欲保留的核心規劃原則或資料流不變條件。 \\[0.35em]
        \texttt{key\_patterns} &
        此經驗所包含的關鍵資料流或規劃模式，用於提醒代理人應保留的中間集合、驗證條件或操作順序。 \\[0.35em]
        \texttt{how\_to\_apply} &
        % TODO-FINAL-VERSION: how_to_apply 已於 v4 補入，但其是否用於索引、直接提示注入，或 v2 的適用性評分／約束化流程，需依最終版本於第四章寫清楚。
        將此經驗套用到新任務時可參考的抽象步驟，但不對應到特定 API 呼叫序列。 \\[0.35em]
        \texttt{pitfalls} &
        此經驗欲避免的常見錯誤，例如過早執行寫入、忽略篩選條件、未完整取得資料或使用錯誤目標集合。 \\
        \hline
    \end{tabularx}
\end{table}

上述欄位使原子化經驗同時具備可追溯性、可檢索性與可注入性。其中，來源任務資訊主要用於建立檢索語境；適用情境、核心原則、關鍵模式與常見陷阱則用於幫助代理人在新任務中辨識相似規劃結構。為避免經驗過度貼近單一訓練任務，本研究在經驗表示上避免保留具體 API 名稱、固定參數值或完整任務流程，而以較抽象的資料來源、候選集合、驗證條件、目標集合與輸出契約等概念描述經驗。

\subsection{經驗檢索流程}

CBR-APIKGAgent 透過語意檢索機制，在不同任務階段取得與當前情境相關的原子化經驗。由於初始規劃與錯誤修復所面對的資訊不同，本研究將同一份離線經驗庫建立為兩種檢索視角：規劃導向索引與修復導向索引。

兩種檢索視角的用途如下：

\begin{itemize}
    \item \textbf{規劃導向索引：}用於初始規劃與重新規劃階段，主要以經驗中的來源任務目標、適用情境、核心原則、關鍵模式與套用方式建立比對內容。此索引的目標，是協助系統根據目前任務目標或既有計畫，找到在資料依賴、目標集合建構或操作順序上相似的經驗。
    \item \textbf{修復導向索引：}用於局部修復與重新規劃中的錯誤分析，主要以經驗中的適用情境、核心原則、關鍵模式、套用方式與常見陷阱建立比對內容。此索引的目標，是使系統能根據失敗訊息或錯誤證據，取得與當前錯誤情境相關的修復參考。
\end{itemize}

換言之，檢索時的語意比對包含兩個部分：一方面，系統會依目前所處階段產生查詢內容，例如任務目標、失敗訊息或重新規劃脈絡；另一方面，經驗庫中的每筆原子化經驗會依規劃導向或修復導向索引，選取不同欄位形成可比對的經驗表示。透過此設計，同一筆經驗可依使用情境呈現不同的檢索重點。

依照使用階段不同，經驗檢索流程可分為以下三種情境：

\begin{enumerate}
    \item \textbf{初始規劃階段}

    \begin{itemize}
        \item \textbf{查詢內容：}以任務目標作為主要查詢。
        \item \textbf{檢索視角：}使用規劃導向索引取得候選原子化經驗。
        % TODO-FINAL-VERSION: 「適用性評估」是偏 v2／修正版的概念。若最後採 v4，可弱化為「依語意相關性與任務階段選取經驗」；若採 v2，保留並在第四章補 LLM applicability scoring。
        \item \textbf{適用性評估：}判斷候選經驗是否能對當前任務形成具體且安全的規劃約束，例如資料依賴、目標範圍、候選驗證、操作順序、批次完整性或輸出限制等。
        % TODO-FINAL-VERSION: 「通過評估、轉換為規劃約束、檢查計畫是否反映約束」偏 v2。若最後採 v4，改成「檢索經驗會被整理並注入提示，作為規劃時的策略參考」。
        \item \textbf{使用方式：}通過評估的經驗會被轉換為規劃約束，注入初始規劃提示中，並用於輔助檢查產生之計畫是否反映相關約束。
    \end{itemize}

    \item \textbf{局部修復階段}

    \begin{itemize}
        \item \textbf{查詢內容：}以失敗訊息、失敗步驟與當前錯誤脈絡作為查詢。
        \item \textbf{檢索視角：}使用修復導向索引取得與錯誤情境相關的候選經驗。
        % TODO-FINAL-VERSION: 「適用性評估」偏 v2／修正版。若最後採 v4，可改為「依錯誤情境檢索並提供相關修復經驗」。
        \item \textbf{適用性評估：}判斷經驗是否有助於解釋當前失敗或形成具體修復依據，例如補足前置資料、修正操作目標、維持目標來源、補齊部分執行、避免副作用或處理其他與錯誤脈絡相關的修復需求。
        % TODO-FINAL-VERSION: 「通過評估」偏 v2；v4 可保留「檢索經驗被整理為修復引導」但不強調評分門檻。
        \item \textbf{使用方式：}通過評估的經驗會被整理為修復引導，供局部修復流程產生插入、替換或刪除步驟時參考。
    \end{itemize}

    \item \textbf{重新規劃階段}

    \begin{itemize}
        \item \textbf{查詢內容：}同時考量任務目標、既有計畫、失敗步驟與錯誤原因。
        \item \textbf{檢索視角：}同時使用規劃導向與修復導向索引，分別取得與任務結構及失敗情境相關的候選經驗。
        % TODO-FINAL-VERSION: 重新規劃階段的「分別評估」偏 v2／修正版。若最後採 v4，改為「分別取得並整理」即可。
        \item \textbf{適用性評估：}分別評估規劃經驗與修復經驗是否適用於當前重新規劃情境。
        % TODO-FINAL-VERSION: 「規劃約束與修復引導」可支援 v2；若 v4 final，建議弱化為「策略參考與修復提醒」。
        \item \textbf{使用方式：}通過評估的經驗會被整合為重新規劃時的規劃約束與修復引導，使新計畫能同時回應原始任務需求與已觀察到的錯誤原因。
    \end{itemize}
\end{enumerate}

% TODO-FINAL-VERSION: 本段最後一句提到「適用性評估具體評分方式、候選經驗篩選門檻、約束檢查」，偏 v2／修正版。若 final 採 v4，需改為 direct retrieval + prompt injection + how_to_apply；若採 v2，第四章補 scoring、selected experiences、constraints 與檢查流程。
整體而言，此檢索流程的重點在於讓不同階段以不同角度使用同一批離線經驗。初始規劃偏重「任務結構是否相似」，局部修復偏重「錯誤情境是否相似」，重新規劃則同時參考規劃導向與修復導向經驗，以協助代理人在修正後續計畫時兼顧原任務目標與已觀察到的失敗原因。原子化經驗重用機制提供的是可檢索並可轉換為規劃約束或修復引導的案例式知識，而非固定規則或保證正確的修復模板；適用性評估的具體評分方式、候選經驗篩選門檻、提示注入位置，以及計畫產生後如何檢查經驗約束是否被反映，則於第四章系統實作中進一步說明。

\section{動態局部修復機制}

動態局部修復機制的目的，是在任務執行過程中發生局部失敗時，先嘗試針對失敗步驟附近的計畫進行小範圍調整，而不是立即重新產生整體計畫。此設計的核心假設是：部分 API 任務失敗並非源自整體任務理解錯誤，而可能是缺少前置資料、使用不適合的 API、資料依賴未正確銜接，或某一步驟的操作目標不完整所造成。若此類問題能透過局部插入、替換或刪除步驟處理，便有機會保留原計畫中已正確完成的部分，並減少重新規劃對既有資料流程造成干擾的可能性。

在本研究所探討的多步驟 API 任務中，代理人執行過程可能已對環境狀態造成改變。因此，錯誤處理流程並不假設系統能將環境回復至任務初始狀態或錯誤發生前狀態。無論是既有重新規劃機制或本研究提出的局部修復機制，皆是在目前已觀察到的執行狀態下產生後續行動；兩者差異在於，重新規劃傾向重新產生較大範圍的後續計畫，而局部修復則優先嘗試對失敗步驟附近進行小範圍調整。

基於上述設定，本研究不以可驗證的錯誤根因診斷作為局部修復之前提。若要精確驗證不同錯誤根因假設，通常需要能夠回復環境狀態、重放特定步驟，或比較多個修復分支；然而，此類流程較接近離線除錯或搜尋式求解，並非本研究所關注的單次任務執行脈絡下之代理人錯誤恢復能力。需要注意的是，局部修復並不保證每次皆能改善任務結果。本研究將局部修復視為一種受限制的錯誤恢復嘗試：系統根據目前可取得的錯誤訊息、失敗步驟、相依資料與案例式修復經驗，判斷是否存在可行的局部調整方案；若證據不足或修復嘗試超出可控範圍，則回退至重新規劃流程。

\subsection{錯誤資訊與相依資料蒐集}

局部修復的品質高度依賴失敗發生時可取得的上下文資訊。若系統只根據單一錯誤訊息選擇修復 API，容易忽略原任務目標、前後步驟關係或已完成資料流程，導致修復方向偏離原計畫。因此，本研究在步驟失敗後，會先於任務執行器端整理錯誤訊息並產生暫定錯誤分類；若錯誤處理路由判斷可進入局部修復，修復流程再進一步使用這些資訊搜尋候選 API、檢索案例式修復參考並產生局部計畫修改。

依照實際處理順序，錯誤資訊與相依資料蒐集可分為下列四個步驟：

\begin{enumerate}
    \item \textbf{記錄失敗步驟與原始錯誤訊息：}當某一步驟執行失敗時，系統會保留該步驟在整體計畫中的位置、步驟描述、原本欲完成的局部目標，以及 API 呼叫或資料處理所回傳的錯誤訊息。
    \item \textbf{整理 API 回饋與相依資料脈絡：}系統會進一步保存失敗 API 的相關資訊、錯誤細節、前序步驟產生的中間結果，以及目前可觀察到的資料依賴脈絡。這些資訊用於判斷失敗是否可能與缺少前置資料、資料來源不一致或操作目標不完整有關。
    \item \textbf{產生暫定錯誤分類與可行回饋：}在錯誤處理路由做出決策前，系統會參考 AgentErrorTaxonomy \cite{zhu2025agentdebug} 所整理之代理人錯誤分類概念，並根據失敗訊息、失敗步驟、API 回饋與相依資料，由大型語言模型產生兩項核心輸出：

    \begin{itemize}
        \item \textbf{暫定錯誤分類（\texttt{error\_type}）：}描述目前失敗較可能接近何種錯誤情境，例如參數缺漏、資料查詢失敗、權限或狀態不符，或原步驟可能需要拆分與較大範圍調整等。在本研究實作中，若原始失敗步驟被暫定分類為不可行計畫，錯誤處理路由會直接回退至重新規劃。
        \item \textbf{可行回饋（\texttt{actionable\_feedback}）：}將錯誤訊息整理為後續修復可使用的操作建議，例如補足前置資料、重新搜尋候選 API，或避免重複執行已完成動作。此回饋不直接決定是否進入局部修復；若系統進入局部修復流程，則會作為修復 API 搜尋與修復策略選擇的重要依據。
    \end{itemize}

    \item \textbf{交由錯誤處理路由與局部修復流程使用：}錯誤處理路由會先根據暫定錯誤分類與修復限制判斷是否進入局部修復。若進入局部修復，上述錯誤證據與可行回饋會被用於修復 API 搜尋；後續修復策略選擇時，系統再結合候選 API、相依資料與案例式修復參考，判斷應插入、替換或刪除局部計畫步驟。
\end{enumerate}

需要說明的是，本研究雖參考錯誤分類研究整理失敗資訊，但此錯誤情境判斷並不等同於可驗證的錯誤根因診斷，也不代表系統能正確辨識真正的上游錯誤原因。多步驟 API 任務中的錯誤原因往往需要結合失敗訊息、原步驟目的、前序資料與候選 API 才能推估；同一個執行錯誤也可能來自不同上游原因。因此，本研究將錯誤情境判斷視為修復流程中的輔助訊號，而非決定修復結果的唯一依據。後續修復策略選擇時，系統會再次結合候選 API、案例式修復經驗與相依資料進行判斷，而不是只依據初步錯誤分類直接決定修復方式。

\subsection{錯誤處理路由與局部修復觸發條件}

在任務執行器完成錯誤資訊整理後，系統會透過錯誤處理路由判斷是否進入局部修復流程。因此，局部修復並非由多種獨立事件任意觸發，而是以「步驟執行失敗」作為主要觸發點，再根據暫定錯誤分類、失敗步驟類型、修復嘗試次數與修復鏈深度決定後續方向。

在本研究架構中，錯誤處理路由會先判斷目前失敗步驟是否屬於修復流程新增或替換出的步驟。之所以需要區分這兩類情境，是因為修復步驟再次失敗時，不應被視為新的獨立失敗點，而應回掛至其原本欲修復的原始失敗步驟，並沿用同一失敗點的修復嘗試次數與修復鏈深度限制。依此路由規則，進入局部修復的情境可整理為兩類：

\begin{itemize}
    \item \textbf{原始計畫步驟失敗：}當原始計畫中的 API 呼叫或資料處理步驟回傳錯誤，或其結果不足以支撐後續計畫時，若該步驟未被暫定分類為不可行計畫，且該步驟尚未達到局部修復嘗試上限，系統會優先嘗試局部修復。
    \item \textbf{修復步驟再次失敗：}若先前插入或替換出的修復步驟執行失敗，系統會透過修復步驟與原始步驟之間的對應關係，回到其欲修復的原始失敗點，檢查該原始失敗點是否仍有修復嘗試額度與修復鏈深度空間；若尚未超出限制，則再次嘗試產生新的局部修復方案。
\end{itemize}

上述兩類情境皆以「同一失敗點仍在可控修復範圍內」為前提。若原始失敗步驟被暫定分類為不可行計畫，或同一原始失敗點已達修復嘗試上限與修復鏈深度限制，系統即不再進入局部修復，而是回退至重新規劃。換言之，錯誤處理路由主要依據失敗步驟是否為修復步驟、暫定錯誤分類中的不可行計畫訊號，以及局部修復嘗試限制，決定後續應進入局部修復或重新規劃；可行回饋不直接決定路由，而是在進入局部修復後作為修復行動的依據。

\subsection{修復 API 搜尋}

在取得錯誤證據與可行回饋後，系統會重新搜尋可能用於修復的候選 API。與初始規劃階段以使用者任務為主要查詢不同，修復 API 搜尋的查詢重點更接近「如何讓目前失敗步驟恢復可執行狀態」或「是否有其他 API 能完成失敗步驟原本應達成的局部目的」。

本研究將修復階段使用的資訊分為「搜尋查詢」與「後續判斷上下文」兩個層次：

\begin{itemize}
    \item \textbf{搜尋查詢：}主要由失敗訊息與可行回饋組成。失敗訊息提供目前執行失敗的表面徵象，可行回饋則將錯誤資訊轉換為較具操作方向的修復需求，例如補足缺漏資料、重新取得目標物件，或避免不適當的 API 呼叫。
    \item \textbf{後續判斷上下文：}包含失敗步驟內容與其局部目標、候選 API 說明、結構化失敗證據，以及修復導向經驗。其中，結構化失敗證據可包含錯誤代碼、錯誤細節、失敗 API 名稱、結果分析、相依資料脈絡與 API 呼叫經驗等資訊。這些資訊主要用於後續修復策略選擇，而不是單純全部併入 API 搜尋查詢中。
\end{itemize}

透過上述區分，搜尋階段可聚焦於取得可能有用的工具；在取得候選 API 後，系統會進一步檢索修復導向經驗，並將候選 API、結構化失敗資訊與修復經驗一併提供給局部修復模組，用以判斷候選 API 是否適合插入、替換或刪除原計畫步驟。

修復 API 搜尋並不代表系統必然會加入新的 API。候選 API 取得後，局部修復模組仍需在同一個策略選擇過程中，根據錯誤證據、失敗步驟內容、候選 API、相依資料脈絡與修復經驗，判斷候選 API 與原失敗步驟之間的關係。舉例而言，候選 API 可能只是用來補足原步驟缺少的前置資料，也可能直接承接原步驟欲完成的操作；在部分情境下，錯誤證據亦可能顯示原步驟本身不應繼續執行。此中間判斷的目的，是降低後續選擇插入、替換或刪除策略時的模糊性，而非建立完整的修復意圖分類。

\subsection{局部計畫修復策略}

本研究採用插入步驟（Insert）、替換步驟（Replace）與刪除步驟（Delete）三種局部計畫修復策略。三種策略皆以最小化計畫變動為原則，僅調整失敗步驟附近的計畫片段，並盡可能保留已成功執行的步驟與可用中間資料。實際策略選擇時，系統會將錯誤證據、候選修復 API、結構化失敗資訊與修復導向經驗提供給大型語言模型，由其判斷修復步驟與原失敗步驟之間的關係，例如是否應保留原失敗步驟繼續執行、是否應以新 API 取代原步驟，或是否應移除原步驟。

\begin{enumerate}
    \item \textbf{插入步驟（Insert）}

    當失敗原因可能來自缺少前置資料或必要條件時，系統會考慮在原失敗步驟前插入新的修復步驟。此策略適用於原步驟本身仍可能正確，但執行前缺少必要輸入的情境。例如，原步驟需要某個物件識別碼，但計畫尚未先查詢符合條件的物件；此時修復模組可插入查詢步驟，再讓原失敗步驟重新執行。

    插入策略的重點在於保留原失敗步驟的角色。新增步驟只負責補足資料或建立前置條件，而不應任意改變原步驟的任務語意。因此，若候選 API 與原步驟具有相同操作目的，而非僅提供前置資料，系統需進一步評估是否應改採替換策略。

    \item \textbf{替換步驟（Replace）}

    當錯誤訊號顯示原失敗步驟本身可能使用了不適當的 API、錯誤的操作方式，或其描述無法完成原本的局部目標時，系統會考慮以新的修復步驟替換原步驟。此策略適用於候選 API 能直接承接原步驟目的，且保留原步驟反而可能造成重複執行或錯誤副作用的情境。

    替換策略比插入策略更可能影響資料流程，因此本研究採取較保守的使用原則。系統需根據失敗證據、候選 API 功能與相依資料，確認新步驟能合理承接原步驟在計畫中的位置；若候選 API 僅能提供輔助資料，則較適合使用插入策略，而非直接取代原步驟。

    \item \textbf{刪除步驟（Delete）}

    當失敗步驟被判斷為冗餘、重複或不應繼續執行時，系統會考慮將其從局部計畫中移除。此策略適用於原步驟與目前任務狀態不一致，或該步驟可能重複處理已完成操作的情境。

    刪除策略的風險在於可能移除後續步驟仍需要的資料來源，因此系統在刪除步驟後，仍需維持計畫中步驟順序與相依關係的一致性。若刪除後導致後續計畫缺少必要資料，則局部修復不應被視為可行方案。
\end{enumerate}

無論採用何種策略，局部修復完成後，系統皆需更新局部計畫，使後續執行流程能接續修復後的步驟順序與資料依賴。此處的更新包含調整步驟位置、維持或重建相依關係，以及確認修復步驟與原任務目標仍保持一致。實際計畫修改、步驟索引調整與狀態追蹤方式，將於第四章系統實作中說明。

\subsection{局部修復與重新規劃的切換條件}

局部修復與重新規劃並非互相排斥的機制，而是錯誤處理流程中的兩個層次。局部修復負責處理可由小範圍計畫調整嘗試恢復的失敗；重新規劃則作為回退機制，用於處理局部修復無法形成可行方案或修復嘗試失敗的情境。

本研究將下列情況視為切換至重新規劃的依據：

\begin{itemize}
    \item \textbf{原始步驟被判斷為不可行計畫：}當原始失敗步驟被暫定分類為不可行計畫時，系統不嘗試局部修復，而是直接回退至重新規劃。
    \item \textbf{缺乏可用修復 API 或刪除依據：}若候選 API 無法補足前置條件、替代原步驟或完成缺漏操作，且失敗步驟也不適合直接刪除，局部修復流程無法產生計畫變更。
    \item \textbf{修復嘗試達到限制：}同一失敗位置經過有限次局部修復後仍無法恢復執行，為避免重複嘗試造成迴圈，系統停止局部修復。
    \item \textbf{修復鏈深度達到限制：}若修復步驟本身再次失敗，且同一失敗點已累積過多層修復步驟，系統停止繼續擴張局部修復鏈。
\end{itemize}

透過上述切換條件，CBR-APIKGAgent 採取「優先嘗試局部修復，但保留重新規劃回退」的錯誤處理策略。此設計一方面使系統能在局部錯誤發生時嘗試保留既有正確步驟，另一方面也避免在錯誤原因不明或修復方向不可靠時持續修改同一段計畫。第五章將進一步透過實驗與案例分析，評估此局部修復機制對任務完成能力與錯誤修復表現的影響。

\end{ZhChapter}
